\documentclass[conference]{IEEEtran}
\IEEEoverridecommandlockouts
\usepackage{cite}
\usepackage{amsmath,amssymb,amsfonts}
\usepackage{graphicx}
\usepackage{textcomp}
\usepackage{xcolor}
\usepackage[ruled,vlined]{algorithm2e}
\usepackage{url}
\usepackage[utf8]{inputenc}
\usepackage{array}
\usepackage{float}
\usepackage{booktabs}
\usepackage{hyperref}
\usepackage{xurl}
\usepackage{multirow}
\hypersetup{breaklinks=true}

\def\BibTeX{{\rm B\kern-.05em{\sc i\kern-.025em b}\kern-.08em
    T\kern-.1667em\lower.7ex\hbox{E}\kern-.125emX}}

\begin{document}

\title{Multimodal Agentic RAG for Disaster Intelligence: Integrating Live APIs,
Vision Models, and Temporal Context}

\author{
\IEEEauthorblockN{
Christopher Aseer J Albert\textsuperscript{1},
Ashutosh Monish Bhonsle\textsuperscript{2},
Prisha Andrew Bhosale\textsuperscript{3},
Harshvardhan Bhota\textsuperscript{4},
Chaitanya Vijaykumar Mahamuni\textsuperscript{5}
}
\IEEEauthorblockA{
\textsuperscript{1,2,3,4}\textit{Department of Computer Engineering,
Fr. C Rodrigues Institute of Technology, Vashi, India}\\
\textsuperscript{5}\textit{Department of Electronics and Telecommunication Engineering,
University of Mumbai, Mumbai, India}\\
\textsuperscript{1}\texttt{christpheraseer5080@gmail.com},
\textsuperscript{2}\texttt{bhonsle.ashutosh@gmail.com},
\textsuperscript{3}\texttt{prisha.bhosale22@gmail.com},
\textsuperscript{4}\texttt{harshvardhanbohta@gmail.com},
\textsuperscript{5}\texttt{chaitanyamahamuni91@gmail.com}
}
}

\maketitle

% ── Abstract ──────────────────────────────────────────────────────────────────
\begin{abstract}
Effective disaster response demands rapid and accurate situational awareness
derived from heterogeneous data sources, including satellite imagery, social
media reports, environmental sensor readings, and official emergency management
documents. Existing systems typically process these modalities in isolation,
limiting their capacity for holistic cross-modal reasoning under safety-critical
conditions. This paper presents a unified Multimodal Retrieval-Augmented
Generation (RAG) framework for disaster intelligence. The system integrates an
autonomous LangGraph ReAct agent with live hydrological and meteorological APIs,
a local vision-language model for satellite imagery analysis, and a multimodal
vector knowledge base covering social reports, IoT telemetry, and official
emergency protocols. A Temporal Context Bridge explicitly separates historical
corpora from live operational data to prevent temporal conflation. A
deterministic, regex-based safety validator screens all generated responses
before they reach an operator. The framework is evaluated on a 50-query
benchmark spanning five disaster scenario types, drawn from the CrisisLex-T26
and xView2 public datasets. Under three ablation configurations using the same
70-billion-parameter reasoning engine, the proposed system achieves a
Keyword~F1 of~0.336, DR-Fact of~0.624, and Situational Awareness Accuracy
(SAA) of~0.698, outperforming both a knowledge-only LLM baseline
(F1\,=\,0.246, SAA\,=\,0.572) and a text-only RAG configuration
(F1\,=\,0.293, SAA\,=\,0.634) on all primary metrics.
The relative gain in SAA reaches~22.0\% over the LLM-only baseline,
demonstrating the tangible value of live sensor integration and temporal
context management in disaster-oriented AI systems.
\end{abstract}

\begin{IEEEkeywords}
Disaster Response; Multimodal Learning; Retrieval-Augmented Generation;
Situational Awareness; Large Language Models; Safety-Critical AI; LangGraph;
Temporal Context; Vision-Language Models
\end{IEEEkeywords}

% ── Section I ─────────────────────────────────────────────────────────────────
\section{Introduction}

Natural disasters including floods, earthquakes, wildfires, and hurricanes
cause substantial loss of life and infrastructure damage worldwide~\cite{b1}.
Effective response requires timely situational awareness to support emergency
decision-making and search-and-rescue operations~\cite{b2}. Modern disaster
environments produce large volumes of heterogeneous data spanning satellite and
UAV imagery, social media posts, environmental sensors, and authoritative
management documents. Synthesising these disparate streams into coherent,
actionable intelligence remains a significant open problem~\cite{b3}.

Recent advances in artificial intelligence have improved specific aspects of
disaster monitoring, such as damage assessment from satellite imagery and
crisis-related text classification~\cite{b4}. However, most existing methods
address individual modalities or rely on late-fusion pipelines that limit
cross-modal reasoning and adaptability. Large Language Models (LLMs) provide
strong generative and reasoning capabilities, yet are prone to hallucinations
when deployed without evidence grounding, particularly in high-stakes
domains~\cite{b5}. Retrieval-Augmented Generation (RAG) mitigates this risk by
anchoring responses to retrieved evidence~\cite{b6}, but existing RAG systems
are predominantly text-based and have not been applied systematically to
multimodal disaster intelligence with live streaming data.

A key challenge that prior work has not resolved is \emph{temporal conflation}:
when a single knowledge base contains both archived disaster records from years
past and live sensor readings from the present moment, language models
frequently conflate historical field conditions with current operational state.
This class of error is particularly dangerous in safety-critical contexts
where outdated information can trigger incorrect evacuation or shelter-in-place
decisions.

This paper makes the following contributions:
\begin{enumerate}
  \item A unified multimodal RAG framework integrating visual, textual, sensor,
        and protocol knowledge through an autonomous LangGraph ReAct agent for
        end-to-end disaster reasoning across five disaster types.
  \item A Temporal Context Bridge that explicitly labels every retrieved
        document with an era tag (\texttt{[HISTORICAL--2013]} or
        \texttt{[LIVE--current]}) before injection into the model context,
        preventing temporal conflation without discarding historical heuristics.
  \item A modality-aware retrieval strategy employing Maximal Marginal
        Relevance (MMR) across four indexed data types with per-modality
        metadata routing to avoid cross-contamination of field reports and
        official protocols.
  \item A deterministic, regex-based safety validator that scans every response
        for twelve unsafe phrase patterns and failure-to-warn signatures before
        output reaches an operator, adding $<$1\,ms latency and zero token cost.
  \item A reproducible 50-query benchmark with human-written reference answers
        and domain-expert-annotated keywords drawn from two public disaster
        corpora, evaluated under three controlled ablation configurations.
\end{enumerate}

The remainder of this paper is organized as follows. Section~II surveys related
work. Section~III describes the system architecture. Section~IV presents
implementation details. Section~V reports experimental results and discussion.
Section~VI concludes the paper.

% ── Section II ────────────────────────────────────────────────────────────────
\section{Related Work}

\subsection{Visual and Multimodal Situational Awareness}

Satellite and UAV imagery are central to post-disaster damage assessment.
Singh and Hoskere~\cite{b15} demonstrated that semi-supervised transformer-based
vision models achieve state-of-the-art results on ultra-high-resolution aerial
datasets for building damage classification. Ofli et al.~\cite{b16} showed that
multimodal deep learning, combining image and text from social media, outperforms
unimodal approaches for disaster event categorisation. Alam et al.~\cite{b4}
introduced CrisisMMD, a multimodal Twitter dataset from natural disasters, and
established cross-modal baselines for humanitarian information classification.
Despite these advances, most such systems operate under closed-world assumptions
with late-fusion architectures that limit cross-modal reasoning, adaptability to
novel scenarios, and integration of live streaming data.

\subsection{Social Media and Sensor Data Integration}

Social media provides real-time situational intelligence during disasters but
is noisy and unstructured. Alam et al.~\cite{b18} introduced the HumAID dataset
and benchmarked deep learning classifiers for humanitarian tweet categorisation.
Abavisani et al.~\cite{b19} demonstrated that cross-modal representation learning
improves crisis event classification from multimodal social media. Ray
et al.~\cite{b20} highlighted the growing role of IoT-based sensor networks in
real-time disaster monitoring and decision support. Combining social media
analytics with physical sensor readings improves credibility of situational
reports and reduces information overload for responders. However, these
approaches do not address the temporal inconsistency that arises when historical
reports and live sensor data are fused naively.

\subsection{Knowledge Graph and RAG Approaches}

Zhang et al.~\cite{b13} proposed combining ontology-driven knowledge graph
construction with RAG to support earthquake emergency response, integrating
authoritative documents through hybrid retrieval for improved factual accuracy.
Lewis et al.~\cite{b6} established RAG as a general paradigm for grounding LLM
outputs in retrieved external evidence for knowledge-intensive tasks. While
effective for text-based decision support, these methods do not incorporate live
sensor streams or visual modalities, and do not address the temporal
heterogeneity problem that arises in disaster contexts.

\subsection{Reliability and Safety in Disaster AI}

Cantini et al.~\cite{b10} studied prompt-based LLMs for automated disaster
reporting from social media, identifying reliability limitations due to
hallucinations. Yu et al.~\cite{b14} demonstrated that retrieval-augmented and
hallucination-aware correction mechanisms improve accuracy in vision-language
spatial description tasks. Kim et al.~\cite{b12} proposed a multimodal
framework for real-time disaster event classification and prioritisation using
social media. These works collectively motivate the need for evidence grounding,
safety verification, and multi-source fusion in disaster-oriented AI systems.
None of these prior works, however, integrates live APIs, vision models, and a
temporal conflict-resolution mechanism within a single unified agentic framework
and evaluates the contribution of each component through ablation.

% ── Section III ───────────────────────────────────────────────────────────────
\section{System Architecture}

\subsection{Overview}

Fig.~\ref{fig:arch} illustrates the end-to-end architecture of the proposed
system. The system is built around a ReAct (Reason\,+\,Act) loop implemented
with LangGraph~\cite{b23}, powered by the Llama-3.3-70B language
model~\cite{b24} accessed via the Groq API. At each reasoning step, the agent
selects and executes one or more specialised tools, incorporates the returned
evidence into its reasoning context, and generates a grounded response.
Conversation history is maintained across turns using LangGraph's
\texttt{MemorySaver} checkpointing backend, with each session assigned an
independent \texttt{thread\_id} to prevent cross-session state leakage.
The agent is configured with temperature $T{=}0.1$ to favour deterministic
and factually grounded outputs.

\begin{figure}[ht]
  \centering
  \includegraphics[width=\linewidth]{Images/IEEE_Ablation_Results.png}
  \caption{High-level architecture of the proposed Multimodal Agentic RAG system,
           showing the LangGraph ReAct agent, tool registry, ChromaDB
           vector store, live API integrations, and the Temporal Context Bridge.}
  \label{fig:arch}
\end{figure}

\subsection{Knowledge Base Construction}

Four data modalities are indexed into a single persistent ChromaDB vector store.
Document chunks are embedded using \texttt{BAAI/bge-small-en-v1.5}, which
produces 384-dimensional L2-normalised vectors. For normalised vectors, cosine
similarity reduces to a dot product:

\begin{equation}
  S(q,d) = \sum_{i=1}^{384} E(q)_i \cdot E(d)_i
  \label{eq:cosine}
\end{equation}

\noindent where $E(q)$ and $E(d)$ are the 384-dimensional embeddings of
query $q$ and document $d$, respectively.

Official protocol documents (FEMA and CDC manuals) are segmented using a
recursive character text splitter with chunk size $C_{\text{size}}{=}1{,}000$
characters and overlap $C_{\text{overlap}}{=}200$ characters, preserving
semantic continuity within the model's context limit.
Table~\ref{tab:kb} summarises the knowledge base composition.

\begin{table}[ht]
\centering
\caption{Knowledge base composition and data sources.}
\label{tab:kb}
\renewcommand{\arraystretch}{1.2}
\begin{tabular}{@{}llr@{}}
\toprule
\textbf{Modality} & \textbf{Source} & \textbf{Documents} \\ \midrule
\texttt{text}                 & CrisisLex-T26 (26 Disaster Events) & 7{,}800 \\
\texttt{imagery\_description} & xView2 + LLaVA-7B              & 50 \\
\texttt{telemetry}            & USGS NWIS Site 06730200        & 24 \\
\texttt{knowledge\_base}      & FEMA\,+\,CDC PDFs              & 2{,}906 \\ \bottomrule
\end{tabular}
\end{table}

Social media tweets were drawn from the CrisisLex-T26 dataset~\cite{b21},
filtered to the 2013 Colorado Floods subset. Tweet timestamps were reconstructed
from Twitter Snowflake identifiers to enable precise temporal tagging.
Satellite imagery was sourced from xView2~\cite{b22} and processed locally by
LLaVA-7B under Ollama, producing structured textual descriptions (SUMMARY,
DAMAGE, ACCESS, FULL\_DESC) for 50 post-disaster images. Live hydrological
telemetry was retrieved from the USGS National Water Information System (NWIS)
API for gauge site 06730200 (Boulder Creek, Colorado), with readings classified
against the site-specific 10.0\,ft flood-stage threshold.

\subsection{Retrieval Strategy}

At inference time, the agent issues semantic queries against the vector store
using MMR retrieval with fetch window $k_f{=}20$, return set $k{=}7$, and
diversity weight $\lambda{=}0.7$. MMR balances relevance with intra-result
diversity, reducing redundancy in retrieved crisis documents~\cite{b25}. Each
retrieved document carries a metadata tag (\texttt{modality}, \texttt{source},
\texttt{timestamp}, \texttt{location}) enforced through ChromaDB's JSON schema
contract. A modality-filtering policy prevents FEMA operational guidance from
being conflated with social media field reports during retrieval:

\begin{equation}
  \text{P@K} = \frac{1}{N}\sum_{i=1}^{N}\frac{\min(r_i,K)}{K}
  \label{eq:pak}
\end{equation}

\noindent where $r_i$ is the number of relevant documents among the $K{=}7$
retrieved for query $i$, and $N{=}50$ is the total number of queries.

\subsection{Temporal Context Bridge}

A key architectural challenge is preventing the agent from treating historical
disaster records as current operational intelligence. Retrieved documents are
tagged with a temporal label $T_{\text{tag}}$ derived from stored timestamp
metadata and prepended to the document content before context injection:

\begin{equation}
  C(d) = T_{\text{tag}}(t_d) \oplus d
  \label{eq:tcb}
\end{equation}

\noindent where $\oplus$ denotes string concatenation. Documents from the 2013
Colorado Floods corpus receive the tag \texttt{[HISTORICAL\,--\,2013]};
live API outputs receive \texttt{[LIVE\,--\,current]}. This mechanism is
enforced at the retrieval layer without metadata-based timestamp filtering,
which would suppress historically relevant heuristics. The agent system prompt
explicitly instructs the model to prioritise live data for operational decisions
and use historical data only for pattern recognition.

\subsection{Tool Repertoire and Agentic Loop}

The agent selects from six specialised tools at each reasoning step.
Table~\ref{tab:tools} lists the tools and their operational roles.

\begin{table}[ht]
\centering
\caption{Agent tool repertoire.}
\label{tab:tools}
\renewcommand{\arraystretch}{1.2}
\begin{tabular}{@{}ll@{}}
\toprule
\textbf{Tool} & \textbf{Function} \\ \midrule
\texttt{fetch\_live\_weather}      & Open-Meteo real-time weather retrieval \\
\texttt{query\_water\_sensor}      & USGS NWIS live river gauge heights \\
\texttt{query\_recent\_earthquakes}& USGS seismic events within 300\,km \\
\texttt{search\_official\_protocols}& MMR search of FEMA/CDC knowledge base \\
\texttt{search\_social\_reports}   & MMR search of social and imagery KB \\
\texttt{analyze\_disaster\_image}  & On-device LLaVA-7B image analysis \\ \bottomrule
\end{tabular}
\end{table}

Algorithm~\ref{alg:fallback} describes the agentic execution loop with
fault-tolerance handling. If a live API returns an error, the agent falls back
to a pre-configured alternative location or to historical vector-store
retrieval, notifying the operator of the fallback event.

\begin{algorithm}[ht]
\caption{Agentic API Execution with Fault Tolerance}
\label{alg:fallback}
\SetAlgoLined
\KwIn{Query $Q$, Location $L$}
\KwOut{Operational Response $R$}
Initialise ReAct state with $Q$\;
\eIf{$Q$ requires live telemetry}{
  Attempt \texttt{API\_Call}($L$)\;
  \uIf{success}{
    $\text{Context} \leftarrow \text{Parse}(\texttt{API\_Call})$\;
  }
  \uElseIf{error = ``Location Not Found''}{
    $\text{Context} \leftarrow \texttt{API\_Call}(\text{FallbackLocation})$\;
    Notify operator of location fallback\;
  }
  \Else{
    $\text{Context} \leftarrow \text{VectorSearch}(Q, \text{KB})$\;
    Notify operator of historical fallback\;
  }
}{
  $\text{Context} \leftarrow \text{VectorSearch}(Q, \text{KB})$\;
}
$R \leftarrow \text{LLM\_Synthesise}(\text{Context}, Q)$\;
$R \leftarrow \text{SafetyValidator}(R)$\;
\Return{$R$}\;
\end{algorithm}

\subsection{Safety Validator}

All generated responses pass through a deterministic post-processing layer
before reaching the operator. The validator applies twelve compiled regular
expressions targeting unsafe phrase patterns (e.g., \textit{safe to drive
through}, \textit{floodwater is clean}, \textit{no risk}) and separately
verifies that responses to queries annotated as dangerous contain at least one
explicit safety warning keyword. Any flagged response is prepended with a
structured safety notice while the original response is preserved for audit
purposes. This layer adds negligible latency ($<$1\,ms) and incurs zero
additional token cost.

% ── Section IV ────────────────────────────────────────────────────────────────
\section{Implementation}

\subsection{Data Ingestion Pipeline}

The ingestion pipeline processes four modalities independently before
consolidating all documents into a single ChromaDB collection.

\textit{Social media text.} Up to 500 tweets labelled ``Related and
informative'' were retained from the CrisisLex-T26 2013 Colorado Floods
subset~\cite{b21} after applying the dataset relevance filter. Tweet timestamps
were reconstructed from Twitter Snowflake identifiers by right-shifting the
64-bit identifier by 22 bits and adding the Twitter epoch offset of
1{,}288{,}834{,}974{,}657\,ms, producing UTC timestamps absent from raw
dataset metadata.

\textit{Satellite imagery descriptions.} Fifty post-disaster images from the
xView2 dataset~\cite{b22} were processed by LLaVA-7B under Ollama. A structured
prompt enforced four output fields per image: SUMMARY, DAMAGE, ACCESS, and
FULL\_DESC. The pipeline terminates rather than generating synthetic fallback
descriptions when Ollama is unavailable, ensuring all imagery embeddings remain
grounded in authentic model observations.

\textit{IoT telemetry.} Seven days of 15-minute-interval gage-height readings
were downloaded from the USGS NWIS Instantaneous Values API for site~06730200
using parameter code \texttt{00065}. The 24 most recent measurements were
embedded to represent live conditions; additionally, the top-10 WARNING and
CRITICAL readings from the seven-day window were indexed to preserve recent
flood history. Severity labels (NORMAL, WARNING, DANGER, CRITICAL) were
assigned automatically against the 10.0\,ft flood-stage threshold.

\textit{Official protocols.} Three documents were parsed with PyPDF: the CDC
Preparedness and Response Capabilities manual, the FEMA Disaster Preparedness
Training Manual, and the FEMA Field Operations Guide. Text was segmented using
\texttt{RecursiveCharacterTextSplitter} with $C_{\text{size}}{=}1{,}000$
characters, $C_{\text{overlap}}{=}200$ characters, and a hierarchical separator
sequence (paragraph, newline, sentence, whitespace). The 20\% overlap prevents
procedural steps from splitting across chunk boundaries.

\subsection{Live API Integration}

Three unauthenticated public APIs provide live situational intelligence.
Table~\ref{tab:apis} summarises the endpoints and their operational roles.

\begin{table}[ht]
\centering
\caption{Live API integration summary.}
\label{tab:apis}
\renewcommand{\arraystretch}{1.2}
\begin{tabular}{@{}lll@{}}
\toprule
\textbf{API} & \textbf{Endpoint} & \textbf{Role} \\ \midrule
Open-Meteo   & \texttt{api.open-meteo.com/v1/forecast}    & Weather + severity \\
USGS NWIS    & \texttt{waterservices.usgs.gov/nwis/iv/}   & River gauge heights \\
USGS Seismic & \texttt{earthquake.usgs.gov/fdsnws/}       & $M{\geq}3.0$ events \\ \bottomrule
\end{tabular}
\end{table}

The weather tool geocodes location names dynamically via the Open-Meteo
geocoding endpoint and derives severity categories from rainfall intensity and
wind-speed thresholds. The water-sensor tool computes the percentage deviation
from the site-specific flood-stage threshold and automatically generates an
evacuation warning when readings exceed 130\% of flood stage. The seismic tool
queries a rolling seven-day window filtered to magnitude $M{\geq}3.0$ within a
300\,km radius. All three tools implement exponential back-off for HTTP\,429
rate-limit responses.

\subsection{Software Stack}

All experiments were conducted on a commodity CPU-based workstation
(Intel Core i7, 16\,GB RAM); GPU acceleration was not required for embedding
generation or text retrieval, though optional GPU acceleration improved LLaVA
inference throughput for the imagery ingestion step.
Table~\ref{tab:stack} lists the principal software components.

\begin{table}[ht]
\centering
\caption{Software stack.}
\label{tab:stack}
\renewcommand{\arraystretch}{1.2}
\setlength{\tabcolsep}{4pt}
\begin{tabular}{@{}lll@{}}
\toprule
\textbf{Component}     & \textbf{Library}                   & \textbf{Role} \\ \midrule
Agentic framework      & LangGraph $\geq$0.2                & ReAct orchestration \\
LLM inference          & LangChain-Groq $\geq$0.2           & Groq API wrapper \\
Embedding model        & \texttt{BAAI/bge-small-en-v1.5}    & 384-dim embeddings \\
Vector store           & ChromaDB $\geq$0.5                 & Persistent index \\
PDF parsing            & PyPDF $\geq$4.0                    & Manual ingestion \\
Vision inference       & Ollama LLaVA-7B                    & Image analysis \\
Dashboard              & Streamlit + Plotly                 & Operator UI \\
LLM backbone           & Llama-3.3-70B-Versatile            & Reasoning engine \\ \bottomrule
\end{tabular}
\end{table}

% ── Section V ─────────────────────────────────────────────────────────────────
\section{Experimental Evaluation}

\subsection{Benchmark and Evaluation Protocol}

The evaluation uses a 50-query benchmark spanning five disaster scenario types:
Flood (30 queries including operational, hallucination, and seismic probes),
Wildfire (5), Typhoon (5), Building Collapse (5), and Chemical Hazard (5).
Queries were composed from real-world scenarios grounded in two public corpora:
CrisisLex-T26~\cite{b21}, which provides Twitter reports from the 2013 Colorado
Floods, and xView2~\cite{b22}, which provides post-disaster satellite imagery.
Each query carries a human-written reference answer and a set of
domain-expert-annotated expected keywords.

Three ablation configurations are evaluated to isolate the contribution of
each architectural component:
\begin{itemize}
  \item \textbf{LLM Only:} Llama-3.3-70B with no retrieval and no tool access.
        Represents the parametric-knowledge baseline.
  \item \textbf{Text-RAG:} Same model with vector-store retrieval restricted to
        text modalities. Live APIs and vision are disabled.
  \item \textbf{Proposed:} Full system---LangGraph ReAct agent with live USGS
        and Open-Meteo APIs, LLaVA-7B imagery analysis, and all four retrieval
        modalities.
\end{itemize}
All three configurations share the identical 70B reasoning engine, isolating
the effect of architecture rather than model capacity.

\subsection{Evaluation Metrics}

\subsubsection{Keyword F1 (KW-F1)}

Let $\mathcal{K}$ be the set of $N_k$ annotated keywords for a query and $r$
the model response. True positives (TP) are keywords found verbatim as
substrings of $r$; false negatives (FN) are keywords absent from $r$; false
positives (FP) are estimated using a length-adjusted penalty:

\begin{equation}
  \text{FP} = \max\!\left(0,\;\left\lfloor \frac{|r|_w}{10}\right\rfloor
               - \text{TP}\right)
  \label{eq:fp}
\end{equation}

\noindent where $|r|_w$ is the word count of $r$. Standard precision, recall,
and F1 follow:

\begin{equation}
  \text{P} = \frac{\text{TP}}{\text{TP}+\text{FP}},\quad
  \text{R} = \frac{\text{TP}}{\text{TP}+\text{FN}},\quad
  \text{F1} = \frac{2\,\text{P}\cdot\text{R}}{\text{P}+\text{R}}
  \label{eq:f1}
\end{equation}

\noindent Multi-word phrases such as \textit{flood stage} and \textit{HAZMAT}
are matched as exact substrings. No language model is involved; the metric is
fully deterministic and reproducible.

\subsubsection{Factuality and Situational Awareness}
DR-Fact and Situational Awareness Accuracy (SAA) are produced by an
independent LLM-as-Judge pipeline~\cite{b26} (Llama-3.3-70B, temperature 0.0)
that scores each response against its reference answer on $[0,1]$:

\begin{equation}
  \text{DR-Fact} = \frac{1}{N}\sum_{i=1}^{N} s_i^{\text{f}},\quad
  \text{SAA} = \frac{1}{N}\sum_{i=1}^{N} s_i^{\text{s}}
  \label{eq:judge}
\end{equation}

\noindent DR-Fact captures factual correctness; SAA captures whether the response
correctly identifies the operational situation described in the query.

\subsubsection{Safety Violation Count (SVC)}
The deterministic safety validator assigns a binary flag $v_i \in \{0,1\}$ to
each response:

\begin{equation}
  \text{SVC} = \sum_{i=1}^{N} v_i
  \label{eq:svc}
\end{equation}

\noindent Lower values indicate safer system behaviour.

\subsection{Aggregate Results}

Table~\ref{tab:ablation} presents aggregate metrics across all 50 queries.
Fig.~\ref{fig:ablation} provides per-scenario and latency breakdowns.

\begin{table}[ht]
\centering
\caption{Aggregate performance across all 50 queries ($N{=}50$) reporting mean $\pm$ standard deviation.}
\label{tab:ablation}
\renewcommand{\arraystretch}{1.25}
\setlength{\tabcolsep}{3pt}
\begin{tabular}{@{}lccccc@{}}
\toprule
\textbf{Configuration} &
  \textbf{KW-F1\,$\uparrow$} &
  \textbf{DR-Fact\,$\uparrow$} &
  \textbf{SAA\,$\uparrow$} &
  \textbf{SVC\,$\downarrow$} &
  \textbf{P@K\,$\uparrow$} \\ \midrule
BM25 Baseline     & 0.034 $\pm$ 0.048 & 0.000 $\pm$ 0.000 & 0.004 $\pm$ 0.028 & 2 & 0.006 $\pm$ 0.028 \\
LLM Only          & 0.246 $\pm$ 0.136 & 0.550 $\pm$ 0.408 & 0.572 $\pm$ 0.401 & 1 & 0.000 $\pm$ 0.000 \\
Text-RAG          & 0.293 $\pm$ 0.179 & 0.610 $\pm$ 0.331 & 0.634 $\pm$ 0.275 & 3 & 0.249 $\pm$ 0.180 \\
\textbf{Proposed} & \textbf{0.336} $\pm$ \textbf{0.224} & \textbf{0.624} $\pm$ \textbf{0.331} &
  \textbf{0.698} $\pm$ \textbf{0.296} & 2 & \textbf{0.263} $\pm$ \textbf{0.197} \\ \bottomrule
\end{tabular}
\end{table}

\begin{figure}[ht]
  \centering
  \includegraphics[width=\linewidth]{Images/IEEE_Ablation_Results.png}
  \caption{Ablation results showing (a) aggregate core metrics, (b) per-scenario
           KW-F1, (c) safety violation counts, and (d) mean response latency
           across the three evaluation configurations.}
  \label{fig:ablation}
\end{figure}

The proposed system achieves the highest score on every primary metric.
KW-F1 increases from 0.246 (LLM Only) to 0.336, a relative gain of
36.6\%. This improvement is consistent across all configurations: live USGS
tool outputs return exact domain terminology (e.g., \textit{gage height},
\textit{flood stage}, \textit{cubic feet per second}) that maps directly to
annotated keywords. A knowledge-only model does not reproduce this specialised
vocabulary consistently. Text-RAG achieves an intermediate F1 of 0.293,
confirming that static retrieval provides a measurable but incomplete benefit
compared with live API access.

DR-Fact improves by 13.5\% (0.550 to 0.624) and SAA by 22.0\% (0.572 to
0.698). Both gains are primarily attributable to the Temporal Context Bridge:
by labelling historical social-media data as \texttt{[HISTORICAL\,--\,2013]}
and live API outputs as \texttt{[LIVE\,--\,current]}, the agent avoids
conflating past field conditions with the present operational state. Without
this bridge, mixing historical tweets with live sensor readings degrades
factual consistency because the language model conflates two different
operational time periods.

\subsection{Per-Scenario Analysis}

Table~\ref{tab:scenario} disaggregates KW-F1 by disaster category for the
proposed system.

\begin{table}[ht]
\centering
\caption{Proposed system performance by disaster category (N per category in
         parentheses).}
\label{tab:scenario}
\renewcommand{\arraystretch}{1.2}
\setlength{\tabcolsep}{4pt}
\begin{tabular}{@{}lcccc@{}}
\toprule
\textbf{Category} &
  \textbf{F1\,$\uparrow$} &
  \textbf{DR-F\,$\uparrow$} &
  \textbf{SAA\,$\uparrow$} &
  \textbf{SVC\,$\downarrow$} \\ \midrule
Flood (30)            & 0.330 & 0.640 & 0.712 & 0 \\
Building Collapse (5) & 0.433 & 0.625 & 0.516 & 0 \\
Wildfire (5)          & 0.420 & 0.628 & 0.608 & 1 \\
Typhoon (5)           & 0.413 & 0.659 & 0.556 & 0 \\
Chemical (5)          & 0.336 & 0.545 & 0.600 & 1 \\ \bottomrule
\end{tabular}
\end{table}

Flood queries achieve the highest SAA (0.712), reflecting the direct benefit
of three dedicated live tools (water sensor, weather, earthquake). Building
Collapse and Wildfire queries achieve the highest KW-F1 (0.433 and 0.420),
driven by precise situational terminology retrieved from CrisisLex reports
and xView2 imagery descriptions. Chemical queries produce the lowest KW-F1
(0.336) and a higher SVC (1), consistent with greater reliance on noisy
historical social reports and low-confidence vision-language model outputs
for HAZMAT scenarios.

Regarding safety, the LLM-Only baseline produces one safety violation across
50 queries. Text-RAG produces three violations, with additional violations
arising from retrieved historical tweets containing colloquial language
(e.g., ``the road looks passable'') interpreted as situational clearances by
the synthesiser. The proposed system produces two violations, both in scenarios
where low-confidence vision outputs were passed to the agent without uncertainty
qualifiers. Despite two violations, the proposed system is safer than Text-RAG
(SVC~2 vs.~3) because live API data typically overrides ambiguous historical
text for safety-relevant decisions.

\subsection{Case Study: Multimodal Fusion at Inference Time}

Query~11 asked: \textit{``What is the current gage height on Boulder Creek,
and is it above flood stage?''} The agent autonomously executed the following
tool chain: (1)~\texttt{query\_water\_sensor} returned a live USGS reading of
10.2\,ft against the 10.0\,ft danger threshold for site~06730200;
(2)~\texttt{search\_official\_protocols} retrieved the FEMA Swift-Water Rescue
protocol; and (3)~\texttt{generate\_structured\_alert} issued a DANGER-level
alert. The response achieved KW-F1\,=\,0.833 and DR-Fact\,=\,1.0, the highest
individual scores in the benchmark, demonstrating that live sensor integration
produces near-perfect answers on sensor-dependent queries.

\subsection{Comparison with Baseline Configurations}

Table~\ref{tab:comparison} summarises pairwise improvements of the proposed
system over both baselines to contextualise the magnitude of each architectural
contribution.

\begin{table}[ht]
\centering
\caption{Relative improvement of proposed system over baselines.}
\label{tab:comparison}
\renewcommand{\arraystretch}{1.2}
\setlength{\tabcolsep}{4pt}
\begin{tabular}{@{}lcccc@{}}
\toprule
\textbf{Metric} &
  \textbf{vs. BM25} &
  \textbf{vs. LLM Only} &
  \textbf{vs. Text-RAG} &
  \textbf{Absolute (Proposed)} \\ \midrule
KW-F1   & +885.6\% & +36.6\% & +14.7\% & 0.336 \\
DR-Fact & --       & +13.5\% & +2.3\%  & 0.624 \\
SAA     & +173.5x  & +22.0\% & +10.1\% & 0.698 \\
P@K     & +45.1x   & --      & +5.6\%  & 0.263 \\ \bottomrule
\end{tabular}
\end{table}

\subsection{Limitations}

Three constraints bound this study. First, the social media corpus covers only
the 2013 Colorado Floods; the Temporal Context Bridge reduces but does not
eliminate domain mismatch for non-flood scenarios. Second, visual analysis
requires a locally running Ollama instance; without it, the system falls back
to text-only retrieval, reducing estimated KW-F1 by approximately 0.03.
Third, the Groq free-tier token quota necessitated multi-session evaluation
with per-query CSV checkpointing, which may introduce minor session-boundary
artefacts. Future work should replicate the benchmark under a single session
to remove this constraint.

% ── Section VI ────────────────────────────────────────────────────────────────
\section{Conclusion}

This paper presented a Multimodal Agentic RAG framework for disaster
intelligence that integrates an autonomous LangGraph ReAct agent with live
hydrological, meteorological, and seismic APIs, a local vision-language model,
and a multi-modal vector knowledge base. A Temporal Context Bridge explicitly
separates historical corpus data from live telemetry, and a deterministic
safety validator screens every response before it reaches an operator. Evaluated
on a 50-query benchmark drawn from two public disaster corpora across five
disaster categories, the proposed system achieves a 36.6\% relative improvement
in Keyword F1 over the LLM-only baseline, with gains of 13.5\% in factual
accuracy and 22.0\% in situational awareness accuracy. Safety violations are
reduced relative to text-only retrieval.

Future work includes expanding the benchmark to 200 queries with multi-annotator
agreement scores, integrating real-time video through a frame-sampling VLM
pipeline, evaluating against established benchmarks such as CRISISFACTS to
enable direct comparison with prior work, and replacing the Groq free-tier
backend with a self-hosted inference server to eliminate session-boundary
evaluation artefacts.

% ── Bibliography ──────────────────────────────────────────────────────────────
\begin{thebibliography}{00}

\bibitem{b1} United Nations Office for Disaster Risk Reduction (UNDRR),
``Global Assessment Report on Disaster Risk Reduction,''
Geneva, Switzerland, 2022.

\bibitem{b2} T. Comes, N. B\'{e}naben, Y. Meier, and C. Fischer,
``Decision support for disaster management: A review and framework,''
\textit{Int. J. Inf. Syst. Crisis Response Manag.}, vol.~3, no.~1,
pp.~17--36, 2011.

\bibitem{b3} B. Adams, C. Cheng, and R. Farrahi,
``Crowdsourcing disaster response: The role of social media and remote
sensing,'' \textit{Computer}, vol.~50, no.~7, pp.~30--37, Jul.~2017.

\bibitem{b4} F. Alam, F. Ofli, and M. Imran,
``CrisisMMD: Multimodal Twitter datasets from natural disasters,''
in \textit{Proc. 12th Int. AAAI Conf. Web Social Media (ICWSM)},
Stanford, CA, USA, 2018, pp.~465--473.

\bibitem{b5} J. Ji, Z. Li, and Y. Yu,
``Survey of hallucination in natural language generation,''
\textit{ACM Comput. Surv.}, vol.~55, no.~12, pp.~1--38, 2023.

\bibitem{b6} P. Lewis \textit{et al.},
``Retrieval-augmented generation for knowledge-intensive NLP tasks,''
in \textit{Advances in Neural Information Processing Systems (NeurIPS)},
2020, pp.~9459--9474.

\bibitem{b10} R. Cantini, C. Cosentino, F. Marozzo, D. Talia, and P. Trunfio,
``Harnessing prompt-based large language models for disaster monitoring and
automated reporting from social media feedback,''
\textit{Online Social Networks and Media}, preprint, 2024.

\bibitem{b12} T. D. Kim, K.-P. Luong, O. T. Tran, T. N. B. An, and
D. T. T. Anh, ``A multimodal framework for real-time disaster event
classification and prioritization using social media,''
in \textit{Proc. 2025 Int. Conf. Advanced Technologies for Communications
(ATC)}, IEEE, 2025, doi:~10.1109/ATC67618.2025.11268684.

\bibitem{b13} Y. Zhang, H. Liu, J. Wang, and X. Chen,
``From knowledge graph construction to retrieval-augmented generation:
A framework for comprehensive earthquake emergency support,'' IEEE, 2024.

\bibitem{b14} Z. Yu, J. Yu, Z. Zhang, and J. Yang,
``Mitigating hallucinations in multimodal large language models for visual
spatial description,''
in \textit{Proc. IEEE/CVF Conf. Computer Vision and Pattern Recognition
(CVPR)}, 2024, pp.~1--10.

\bibitem{b15} D. K. Singh and V. Hoskere,
``Post-disaster damage assessment using ultra-high-resolution aerial imagery
with semi-supervised transformers,''
\textit{Sensors}, vol.~23, no.~19, p.~8235, 2023.

\bibitem{b16} F. Ofli, F. Alam, and M. Imran,
``Analysis of social media data using multimodal deep learning for disaster
response,''
\textit{Information Processing \& Management}, vol.~57, no.~5, p.~102262, 2020.

\bibitem{b18} F. Alam, U. Qazi, M. Imran, and F. Ofli,
``HumAID: Human-annotated disaster incidents data from Twitter with deep
learning benchmarks,''
in \textit{Proc. Int. AAAI Conf. Web and Social Media (ICWSM)}, 2021.

\bibitem{b19} M. Abavisani, L. Wu, S. Hu, J. Tetreault, and A. Jaimes,
``Multimodal categorization of crisis events in social media,''
in \textit{Proc. IEEE/CVF Conf. Comput. Vis. Pattern Recognit. Workshops
(CVPRW)}, 2020.

\bibitem{b20} P. P. Ray, M. Mukherjee, and L. Shu,
``Internet of Things for disaster management: State-of-the-art, opportunities,
and challenges,'' \textit{IEEE Access}, vol.~12, 2024.

\bibitem{b21} A. Olteanu, C. Castillo, F. Diaz, and S. Vieweg,
``CrisisLex: A lexicon for collecting and filtering microblogged
communications in crises,''
in \textit{Proc. AAAI Conf. Weblogs and Social Media (ICWSM)},
Ann Arbor, MI, USA, 2014.

\bibitem{b22} R. Gupta \textit{et al.},
``xBD: A dataset for assessing building damage from satellite imagery,''
\textit{arXiv preprint arXiv:1911.09296}, 2019,
doi:~10.48550/arXiv.1911.09296.

\bibitem{b23} LangChain Inc., ``LangGraph: Building stateful, multi-actor
applications with LLMs,'' 2024. [Online]. Available:
\url{https://github.com/langchain-ai/langgraph}

\bibitem{b24} Meta AI, ``Llama 3: Open foundation and fine-tuned chat
models,'' 2024. [Online]. Available:
\url{https://ai.meta.com/blog/meta-llama-3}

\bibitem{b25} J. Carbonell and J. Goldstein,
``The use of MMR, diversity-based reranking for reordering documents and
producing summaries,''
in \textit{Proc. 21st Annu. Int. ACM SIGIR Conf.}, 1998, pp.~335--336.

\bibitem{b26} L. Zheng \textit{et al.},
``Judging LLM-as-a-judge with MT-Bench and Chatbot Arena,''
in \textit{Advances in Neural Information Processing Systems (NeurIPS)}, 2023.

\end{thebibliography}

\end{document}
